%% kapa: exam=1 round=26 year=2015 period=1 form=A subject=civil_law subject_ko=민법 seq=6 number=6
%% source: Q-Net 감정평가사 기출문제 (cst00309) artlSeq=5138001 file=extracted/1008/5138001/제26회_감정평가사 1차_기출/제26회_감정평가사 1차_기출/1교시_감평_A형.pdf
%% conversion: kordoc 4.24.51 -> tools/gen1.py (auto); review=proofread-sonnet; flags=none
\begin{question}{6}
권리의 객체에 관한 설명으로 옳지 않은 것은? (다툼이 있으면 판례에 따름)
\begin{choices}
  \item 주물과 종물의 법리는 물건 상호간에 적용되고, 권리 상호간에는 적용되지 않는다.
  \item 천연과실은 그 원물로부터 분리하는 때에 이를 수 취할 권리자에게 속한다.
  \item 독립된 부동산으로서의 건물이라고 하기 위하여는 최소한의 기둥과 지붕 그리고 주벽이 갖추어져야 한다.
  \item 종물은 주물의 처분에 따른다는 민법규정은 임의 규정이다.
  \item 명인방법을 갖춘 수목의 경우에는 토지와 독립한 거래의 객체가 된다.
\end{choices}
\end{question}
\answer{1}
